Factor off the lowest power of $z$. The remaining polynomial $Q$ has constant and leading coefficients one, degree $d\le5$, and binary coefficients. If all its roots have modulus one, conjugation and reciprocation produce the same multiset of roots, so $Q$ is reciprocal. Conversely, every reciprocal binary polynomial of degree at most five has roots of unity. Here is a complete check. For degrees $0,1,2,3$, they are $1$, $1+z$, $1+z^2$, $1+z+z^2$, $1+z^3$, and $(1+z)(1+z^2)$. For degree four the four possibilities factor as $z^4+1$, $(z^2+z+1)(z^2-z+1)$, $(z+1)^2(z^2-z+1)$, and $1+z+z^2+z^3+z^4$. For degree five the four possibilities factor as $1+z^5$, $(z+1)(z^2+1)(z^2-z+1)$, $(z+1)(z^4+1)$, and $1+z+\cdots+z^5$. Every listed factor has roots of unity, by geometric-series identities or direct quadratic roots. There are respectively $1,1,2,2,4,4$ normalized polynomials of degrees $0$ through $5$. Each degree-$d$ polynomial can be shifted to exponents within $1,\ldots,6$ in $6-d$ ways. The answer is $6+5+8+6+8+4=\boxed{37}$.
